\clearpage\fancyhead[L]{\small GAUSS TO BODE.}

\begin{center}\textbf{GAUSS TO BODE.}\par\vspace{4pt}\hrule\vspace{4pt}\fontsize{10}{12}\selectfont Astronomisches Jahrbuch für 1813. pp. 236\,..\,238. Berlin 1810.\par\vspace{4pt}\hrule\vspace{4pt}\end{center}
\begin{flushright}Göttingen, 2 September 1810.\end{flushright}

\par\indent ---\enspace---\enspace---

\par\indent The elements used to calculate this ephemeris [of Pallas for 16 Sept. 1810 to 1 July 1811 by \textsc{Gauss}] are based only on the last 4 oppositions and will be sufficient to find the planet this time. Some time ago, however, I also calculated the elements that agree as closely as possible with all 6 oppositions observed so far and can now serve as the most suitable basis for calculating the perturbations. The perturbations are now beginning to show up quite clearly. I will shortly submit a memoir on these calculations to our Society of Sciences. You will already know that our government intends to do a great deal for astronomy 200000 francs are to be spent on building the new observatory. Professor \textsc{Harding} has received a grant of 4000 francs for a trip to Paris to fill in the gaps in the Histoire céleste there and complete his star charts.

I greatly value the honor recently shown me by my admission to your Academy of Sciences.
